% BEGIN EMBEDDED BIBLIOGRAPHY
% Bibliography data are part of this source; the compiler writes its temporary input.
\begin{filecontents*}[overwrite]{\jobname.bib}
@unpublished{QED3Master,
  title = {The field-strength {Osterwalder--Schrader} axioms for {QED} in three dimensions},
  date = {2026-10-20},
  note = {Comprehensive edition, 263 pp., 20 October 2026, before the modular integration revision}
}
@article{HeisenbergEuler1936,
  author = {Heisenberg, W. and Euler, H.},
  title = {Folgerungen aus der {Diracschen} Theorie des Positrons},
  journal = {Zeitschrift f{\"u}r Physik},
  volume = {98},
  date = {1936},
  pages = {714--732},
  doi = {10.1007/BF01343663},
  eprint = {physics/0605038},
  eprinttype = {arxiv},
  note = {English translation by W. Korolevski and H. Kleinert}
}
@online{PongeSchatten2025,
  author = {Ponge, Rapha{\"e}l},
  title = {Singular Values and Schatten Classes},
  date = {2025},
  note = {Introduction to Noncommutative Geometry, Chapter 4; slides 16 and 53--56},
  url = {https://raphaelponge.org/wp-content/uploads/2025/07/summer-chap4-1.pdf}
}
@online{GuthFourier2015,
  author = {Guth, Larry},
  title = {18.156 Lecture Notes: {March} 30, 2015},
  date = {2015-03-30},
  organization = {Massachusetts Institute of Technology},
  note = {Transcribed by Jane Wang; Theorems 4--7},
  url = {https://math.mit.edu/~lguth/Math156/notes/20.pdf}
}

@unpublished{ChamberLattice,
 title={Reflection positivity for chamber {QED} on the infinite lattice},
 date={2026-10-20}, note={Companion manuscript, 20 October 2026}}

@unpublished{FermionBlock,
 title={Chamber fermion locality and local integration},
 date={2026-10-20}, note={Companion manuscript, 20 October 2026}}
\end{filecontents*}
% END EMBEDDED BIBLIOGRAPHY
\documentclass[11pt,letterpaper,reqno]{amsart}
% BEGIN INLINE manuscript-style
\makeatletter
\NeedsTeXFormat{LaTeX2e}[2020/10/01]
% Use with \documentclass[11pt,letterpaper,reqno]{amsart}.
% This file controls presentation only. It supplies no mathematical content.

\RequirePackage[T1]{fontenc}
\RequirePackage{lmodern}
\RequirePackage{microtype}
\RequirePackage[letterpaper,margin=1in,headheight=15pt]{geometry}
\RequirePackage{amsmath,amsthm,amssymb,mathtools}
\RequirePackage{enumitem}
\RequirePackage{perpage}
\RequirePackage{xcolor}
\definecolor{manuscriptred}{rgb}{0.6,0.1,0.1}
\definecolor{manuscriptblue}{rgb}{0.1,0.3,0.5}
\definecolor{manuscriptmagenta}{HTML}{60445C}
\RequirePackage[breaklinks=true,colorlinks=true,
  linkcolor=manuscriptred,citecolor=manuscriptblue,urlcolor=manuscriptmagenta]{hyperref}
\RequirePackage{cleveref}
\hypersetup{pdfauthor={}}

% Retain the AMS paragraph, title, heading and running-head conventions.
% Equations retain one global Arabic sequence, as in the reference.
\setcounter{tocdepth}{1}
\renewcommand{\thefigure}{\Roman{figure}}
\renewcommand{\thetable}{\Roman{table}}
\MakePerPage{footnote}
\renewcommand{\thefootnote}{%
  \ifnum\value{footnote}>9\arabic{footnote}\else\fnsymbol{footnote}\fi}

% Upright statements. Main results are A, B, C; supporting statements
% share a section-local sequence. These end marks indicate typography,
% never proof completeness, validation, or a deferred-proof status.
\theoremstyle{definition}
\newtheorem{ms@theorem}{Theorem}
\renewcommand{\thems@theorem}{\Alph{ms@theorem}}
\newtheorem*{ms@theoremstar}{Theorem}
\newtheorem{ms@proposition}{Proposition}[section]
\newtheorem{ms@lemma}[ms@proposition]{Lemma}
\newtheorem{ms@definition}[ms@proposition]{Definition}
\newtheorem{ms@corollary}[ms@proposition]{Corollary}
\newtheorem*{ms@corollarystar}{Corollary}

\newcommand{\ms@majorstart}{%
  \pushQED{\qed}\renewcommand{\qedsymbol}{$\blacksquare\square$}}
\newcommand{\ms@minorstart}{%
  \pushQED{\qed}\renewcommand{\qedsymbol}{$\blacksquare$}}
\newenvironment{theorem}[1][]{\ms@majorstart\ms@theorem[#1]}
  {\popQED\endms@theorem}
\newenvironment{theorem*}[1][]{\ms@majorstart\ms@theoremstar[#1]}
  {\popQED\endms@theoremstar}
\newenvironment{proposition}[1][]{\ms@minorstart\ms@proposition[#1]}
  {\popQED\endms@proposition}
\newenvironment{lemma}[1][]{\ms@minorstart\ms@lemma[#1]}
  {\popQED\endms@lemma}
\newenvironment{definition}[1][]{\ms@minorstart\ms@definition[#1]}
  {\popQED\endms@definition}
\newenvironment{corollary}[1][]{\ms@majorstart\ms@corollary[#1]}
  {\popQED\endms@corollary}
\newenvironment{corollary*}[1][]{\ms@majorstart\ms@corollarystar[#1]}
  {\popQED\endms@corollarystar}
% The source's lemma/proposition variants are deliberately not carried over:
% choose one marker for each statement type rather than encode proof status.

\theoremstyle{remark}
\newtheorem{remark}[ms@proposition]{Remark}
\newtheorem*{remark*}{Remark}
% The ordinary amsthm proof environment retains its open square.

\crefname{ms@theorem}{theorem}{theorems}
\Crefname{ms@theorem}{Theorem}{Theorems}
\crefname{ms@proposition}{proposition}{propositions}
\Crefname{ms@proposition}{Proposition}{Propositions}
\crefname{ms@lemma}{lemma}{lemmas}
\Crefname{ms@lemma}{Lemma}{Lemmas}
\crefname{ms@definition}{definition}{definitions}
\Crefname{ms@definition}{Definition}{Definitions}
\crefname{ms@corollary}{corollary}{corollaries}
\Crefname{ms@corollary}{Corollary}{Corollaries}
\crefname{remark}{remark}{remarks}
\Crefname{remark}{Remark}{Remarks}

% Bibliography configuration for this manuscript.
  \RequirePackage[backend=biber,style=alphabetic,backref=true,
    giveninits=true,isbn=false,date=year]{biblatex}
  \addbibresource{\jobname.bib}
  \setcounter{biburlnumpenalty}{100}
  \setcounter{biburllcpenalty}{100}
  \setcounter{biburlucpenalty}{100}
  \renewcommand*{\bibfont}{\footnotesize}
  \renewbibmacro{in:}{}
  \DeclareFieldFormat{pages}{#1}
  % Use the public alphabetic style. No private author-hash overrides;
  % repeated author names are printed in full.
\makeatother
% END INLINE manuscript-style
\newcommand{\R}{\mathbb R}
\newcommand{\C}{\mathbb C}
\newcommand{\Z}{\mathbb Z}
\newcommand{\N}{\mathbb N}
\newcommand{\dd}{\mathrm d}
\DeclareMathOperator{\supp}{supp}
\DeclareMathOperator{\tr}{tr}
\DeclareMathOperator{\diag}{diag}
\title[Positive probes and free Dirac response]{Positive compact probes and free response coefficients for lattice QED in three dimensions}
\date{}
\hypersetup{pdftitle={Positive compact probes and free response coefficients for lattice QED in three dimensions},pdfauthor={}}
\begin{document}
\begin{abstract}
We construct compactly supported real two-form probes with positive
fourth response for the massive doubled Dirac operator in three
dimensions. The proof gives the local large-mass coefficient
$1/(240\pi M^5)$ and a remainder of order $M^{-7}$.
For the chamber lattice regulator, the fourth logarithmic
determinant derivative converges to this continuum trace with
error $O(\eta^{1/4})$ at fixed physical mass, uniformly in the
torus size; the remaining volume error decays exponentially.
A cancellation between the two-vertex and contact terms gives
uniform bounds on the second and fourth ordinary determinant
derivatives. The constants are uniform under real translations
of the compact probe.
\end{abstract}
\maketitle
\tableofcontents

\section{Introduction}\label{sec:introduction}

Let $\sigma_\nu$ be the Pauli matrices and
$\gamma_\nu=\diag(\sigma_\nu,-\sigma_\nu)$, $1\le\nu\le3$.
Thus $\gamma_\nu\gamma_\mu+\gamma_\mu\gamma_\nu=2\delta_{\mu\nu}$.
On $L^2(\R^3;\C^4)$, for $M>0$ and real
$a\in C_c^\infty(\R^3;\R^3)$, write
\[
 S_M=(M+\gamma\cdot\partial)^{-1},\qquad
 V_a=i\gamma\cdot a,\qquad
 q_{4,M}(a)=-\tfrac14\tr[(S_MV_a)^4].
\]
For a real compact smooth two-form $\beta$, use the Euclidean
codifferential
$(d^*\beta)_\nu=-\sum_{\mu\ne\nu}\partial_\mu\beta_{\mu\nu}$,
put $f=dd^*\beta$, and define, at fixed mass $\bar m>0$,
\begin{equation}\label{eq:positive-coefficient}
 b(f)=-6\tr[(S_{\bar m}i\gamma\cdot d^*\beta)^4]
      =24q_{4,\bar m}(d^*\beta).
\end{equation}
This depends only on $f$: two choices of $d^*\beta$ with the same
exterior derivative differ by a compactly supported smooth one-form
that is closed and coclosed, hence harmonic and therefore zero.
\begin{theorem}[A positive continuum probe]\label{prop:positive-probe}
For every real $a\in C_c^\infty(\R^3;\R^3)$ and $M>0$,
the operator $(S_MV_a)^4$ is trace class and its trace is real. There
is $C_a<\infty$ such that, for $M\ge1$,
\begin{equation}\label{eq:large-mass-quartic}
 \left|q_{4,M}(a)-\frac1{240\pi M^5}
       \int_{\R^3}|\operatorname{curl}a|^4\,\dd x\right|
 \le C_aM^{-7}.
\end{equation}
In particular, for every $\bar m>0$ there is a real
$\beta\in C_c^\infty(\R^3;\Lambda^2)$ for which $f=dd^*\beta$ satisfies
$b(f)>0$.
\end{theorem}


The sign in \eqref{eq:positive-coefficient} is obtained by dilating
the compact probe after evaluating the large-mass coefficient.
The leading local term is the three-dimensional analogue of the
quartic field-strength term in the Heisenberg--Euler effective
Lagrangian \cite{HeisenbergEuler1936}. The coefficient and remainder
in \cref{prop:positive-probe} are proved here by the continuum
argument of \cite[\S17]{QED3Master}.

\subsection{The lattice response}\label{subsec:response-conventions}

Fix $0<\rho<1$ and $0<\bar m\le1$. Let $\beta$ be a real nonzero compactly
supported smooth two-form, and put $a=d^*\beta$, $f=dd^*\beta$.
Thus $a_\nu=-\sum_{\mu\ne\nu}\partial_\mu\beta_{\mu\nu}$.
Let $Q_\beta$ be a closed coordinate cube of least side $D_\beta$
containing $\supp\beta$; this side includes no buffer. For $0<\eta\le1$
define
\[
 J_\eta v(y)=\eta^{-3}\int_{[-\eta/2,\eta/2]^3}v(y+z)\,\dd z,
 \qquad
 \delta_{\eta,\mu}v(y)=
      \frac{v(y+\eta e_\mu/2)-v(y-\eta e_\mu/2)}{\eta}.
\]
The interpolating functions
\begin{equation}\label{eq:response-lattice-probe}
 a_{\eta,\nu}^0(y)=-\sum_{\mu\ne\nu}
               \delta_{\eta,\mu}J_\eta\beta_{\mu\nu}(y)
\end{equation}
are defined without periodization. If $n=R/\eta\in\N$, $n\ge2$,
the lattice one-cochain on the physical torus is
\begin{equation}\label{eq:response-periodic-probe}
 a_{\eta,\nu}(\eta c)=\sum_{z\in\Z^3}
       a_{\eta,\nu}^0(\eta c+\eta e_\nu/2+Rz).
\end{equation}
It is the adjoint lattice coboundary of the cube-averaged probe
at plaquette centers. In lattice units the external connection
is $B_t(c,c+e_\nu)=t\eta a_{\eta,\nu}(\eta c)$, extended with
the orientation sign to reversed bonds. The free mass in
these units is $m=\bar m\eta$.


The chamber labels are the $48$ signed permutation matrices
$g\in O_{\mathrm{cub}}=O(3)\cap\operatorname{GL}(3,\Z)$.
Let $s_0$ change the sign of the third coordinate, let $s_1,s_2$
interchange coordinates $(2,3)$ and $(1,2)$, respectively, and let
$r_{e_1}$ change the sign of the first coordinate. Set
\begin{gather*}
 n_0=e_3,\quad n_1=(e_2-e_3)/\sqrt2,\quad
 n_2=(e_1-e_2)/\sqrt2,\quad n_3=-e_1,\\
 (w_0,w_1,w_2,w_3)=(6,6\sqrt2,6\sqrt2,6).
\end{gather*}
For a spinor $\psi:\Z^3\times O_{\mathrm{cub}}\to\C^4$ with
$\psi(c+ne_\nu,g)=-\psi(c,g)$, put
$M_i(g)=\frac{w_i}{2}(\rho+\gamma\cdot gn_i)$ and define
\begin{equation}\label{eq:prm-chamber-operator}
\begin{split}
 (D^{\mathrm{ch}}[B]\psi)(c,g)
 ={}&\sum_{i=0}^2M_i(g)\bigl(\psi(c,g)-\psi(c,gs_i)\bigr)\\
 &+M_3(g)\bigl(\psi(c,g)
       -e^{iB(c,c+ge_1)}\psi(c+ge_1,gr_{e_1})\bigr).
\end{split}
\end{equation}
Here $B$ is an $n\Z^3$-periodic cochain on the oriented nearest-neighbor
edges of the covering lattice, with $B(c',c)=-B(c,c')$ for
$c,c'\in\Z^3$. The operator is determined by its values for
$c\in\{0,\ldots,n-1\}^3$; when $n=2$, edges with the same quotient
endpoints remain distinct. This is the chamber operator
of \cite[\S2]{ChamberLattice} in lattice units.
The spinor pairing uses counting measure in $c$ and
$48^{-1}\sum_g$ in the labels.

Let
\begin{equation}\label{eq:response-free-ratio}
 N_0(t)=\frac{\det(D^{\mathrm{ch}}[B_t]+m)}
                  {\det(D^{\mathrm{ch}}[0]+m)},\qquad
 c_j=N_0^{(j)}(0),\qquad
 a_{\mathrm{ch}}=(\log N_0)''(0),\quad
 b_{\mathrm{ch}}=(\log N_0)^{(4)}(0).
\end{equation}
The logarithm is the holomorphic branch equal to zero at the
origin. Its denominator is nonzero: the real part of
$D^{\mathrm{ch}}[0]+m$ is at least $m>0$. For comparison
with the charge convention, for $e>0$ put $q=e\eta^{1/2}$; the normalized
probe is $\widehat a=\eta^{1/2}a_\eta$, so
$B_t=q(t/e)\widehat a$. The free ratio and its coefficients
are independent of $e$.

On $T_R=\R^3/R\Z^3$ use spinors antiperiodic in all three directions.
Their momentum set is
$\mathcal P_R=\frac{2\pi}{R}(\Z+\frac12)^3$, and
$S^R$ is the multiplier $s_{\bar m}(P)$ there, where
$s_m(p)=(m+i\gamma\cdot p)^{-1}$. For
$R\ge D_\beta+1$ let $a^R=\sum_za(\,\cdot+Rz)$, and put
\begin{equation}\label{eq:response-continuum-fourth}
 b_R(f)=-6\operatorname{Tr}(S^Ri\gamma\cdot a^R)^4,
 \qquad b(f)=-6\operatorname{Tr}(S_{\bar m}i\gamma\cdot a)^4.
\end{equation}


\begin{theorem}[Free lattice response]\label{prop:response-fourth}
For the operator and sampled probe above, fix
$0<\rho<1$, $0<\bar m\le1$, and real nonzero
$\beta\in C_c^\infty(\R^3;\Lambda^2)$.
For all $0<\eta\le1$ and $n=R/\eta\in\N$ with $n\ge2$,
the following assertions hold.

If $R\ge D_\beta+1$, the operators $S^Ri\gamma\cdot a^R$ and
$S_{\bar m}i\gamma\cdot a$ belong to the Schatten class
of order four; their fourth-power traces are real.
There are constants
$C_1(\rho,\bar m,\beta)$ and $C_2(\bar m,\beta)$ such that
\begin{align}\label{eq:response-fourth-errors}
 |b_{\mathrm{ch}}(\eta,R)-b_R(f)|&\le C_1\eta^{1/4},\notag\\
 |b_R(f)-b(f)|&\le C_2e^{-\bar m(R-D_\beta)},
 \qquad R\ge D_\beta+1,\quad0<\eta\le1,\quad R/\eta\in\N.
\end{align}
The constants are unchanged by every real translate of
the probe. In particular $b_{\mathrm{ch}}(\eta,R)$ tends
to $b(f)$ whenever $\eta\to0$ and $R\to\infty$.

The function $N_0$ is entire, even, and positive on the real axis.
Its derivatives at zero are real and satisfy
\[
 c_{2j+1}=0,\qquad c_2=a_{\mathrm{ch}},\qquad
 c_4=b_{\mathrm{ch}}+3c_2^2 .
\]
There is $C_{\rho,\bar m,\beta}<\infty$, unchanged by real
translations of $\beta$, such that
\[
 |c_2|\le C_{\rho,\bar m,\beta}\quad(R\ge1),\qquad
 |c_4|\le C_{\rho,\bar m,\beta}\quad(R\ge D_\beta+1).
\]
All coefficients are independent of $e$ in the charge convention
above. They are unchanged by translating the probe through
$\eta\Z^3$ or by replacing its external cochain by a cubic reflected
cochain.
\end{theorem}


\subsection{The proof and its use}

For a probe chosen by \cref{prop:positive-probe}, the errors in
\eqref{eq:response-fourth-errors} are eventually smaller than its
positive coefficient. The choice of probe precedes every
interacting comparison. The free estimates themselves apply to
every real compact smooth probe, regardless of the sign of $b(f)$.

The continuum proof differentiates the fourth trace in the mass.
Its diagonal kernel admits a local expansion, and gauge covariance
reduces the leading term to the antisymmetric first derivative of
the connection. A constant-field heat kernel evaluates that term;
integration in the mass and dilation give \cref{prop:positive-probe}.
The kernel regularity and coefficientwise inverse identities are
proved in \cref{sec:positive-probe}.

For the lattice comparison, the chamber inverse splits into a
Dirac inverse on constant labels and a uniformly bounded remainder.
Repeated differentiation of a bond phase gives contact terms.
The contact traces tend to zero, while the four first-order
vertices converge in Schatten norm. Fourier aliases and the
discarded high momenta are bounded uniformly in the torus size.
The comparison with infinite space then follows from the
exponential decay of the Yukawa images. These comparisons are
proved in \cref{subsec:response-fourth}.

The second derivative needs a separate cancellation. Its
two-vertex trace and contact trace form a polarization kernel whose
momentum integral vanishes to second order at zero. Pairing its
quadratic bound and the finite-volume quadrature error with the
decaying probe transform yields the uniform bound for $c_2$.
The identity $c_4=b_{\mathrm{ch}}+3c_2^2$ then gives the bound
for the ordinary fourth derivative. The cancellation is proved in
\cref{subsec:response-polarization}, and
\cref{subsec:response-uniform-coefficients} completes the proof of
\cref{prop:response-fourth}.

The proofs reproduce the continuum and chamber response arguments
of \cite[\S17--18]{QED3Master}. The chamber construction and its
free Schur-complement estimates are imported from that manuscript
with precise references below. The present results concern the
free determinant at zero background connection.

We use Euclidean norms and the Fourier convention
$\widehat v(P)=\int_{\R^3}v(x)e^{-iP\cdot x}\,\dd x$.
Write
$|\beta|_{C^j}=\max_{\mu<\nu,\,|\alpha|\le j}
\|\partial^\alpha\beta_{\mu\nu}\|_\infty$.
For compact operators, $\|\cdot\|_p$ denotes the Schatten norm;
bounded multiplication, completeness, trace cyclicity and
H\"older's inequality are used in the forms of
\cite[slides~16, 53--56]{PongeSchatten2025}.
Finite-dimensional traces include all spin and chamber labels.
\section{The local continuum response}\label{sec:positive-probe}

We prove \cref{prop:positive-probe}. The mass derivative has a
five-propagator kernel whose diagonal is continuous. A compact gauge
transformation removes the constant connection at each point.
Quartic homogeneity and parity then leave only the fourth power of
the field strength at leading order. The constant-field calculation
determines its coefficient.

\subsection{Kernels and differentiation in the mass}

The kernel and Fourier multiplier of $S_M$ are
\begin{align}
 S_M(z)&=(M-\gamma\cdot\partial)
              \frac{e^{-M|z|}}{4\pi|z|},\label{eq:probe-free-kernel}\\
 \widehat S_M(p)&=\frac{M-i\gamma\cdot p}{M^2+|p|^2},\qquad
 |\widehat S_M(p)|=(M^2+|p|^2)^{-1/2}.
 \label{eq:probe-free-symbol}
\end{align}
In particular,
\[
 |S_M(z)|\le\frac{e^{-M|z|}(2M|z|+1)}{4\pi|z|^2},\qquad
 S_M(z)=M^2S_1(Mz),\qquad \|S_M\|_{L^1}\le3/M.
\]
The operator $G=S_MV_aS_M$ is Hilbert--Schmidt. Indeed, its momentum
kernel, with measure $(2\pi)^{-3}\dd p$, is
$\widehat S_M(p)i\gamma\cdot\widehat a(p-q)\widehat S_M(q)$.
Writing $u(p)=(M^2+|p|^2)^{-1}$, the matrix norm bound and Young's
inequality give
\[
 \|G\|_{\mathrm{HS}}^2
 \le8(2\pi)^{-6}\int u(p)(|\widehat a|^2*u)(p)\,\dd p
 \le8(2\pi)^{-3}\|u\|_2^2\|a\|_2^2
 =\frac{\|a\|_2^2}{\pi M}.
\]
Here $\|u\|_2^2=\pi^2/M$. Since $(S_MV_a)^4=GV_aGV_a$,
\begin{equation}\label{eq:probe-trace-bound}
 |q_{4,M}(a)|\le\frac{\|a\|_2^2\|a\|_\infty^2}{4\pi M}.
\end{equation}
The matrix $\gamma_5=\left(\begin{smallmatrix}0&1_2\\1_2&0\end{smallmatrix}\right)$
conjugates $S_M$ to $S_M^*$ and $V_a$ to $V_a^*=-V_a$.
Conjugation of the fourth power and cyclicity of the trace therefore
give a real trace. Cyclicity here involves a bounded factor and the
trace-class product $GV_aG$, so it does not presume that $S_MV_a$
itself is trace class.

We will evaluate the mass derivative on the diagonal of a kernel with
five propagators. This requires continuity at the diagonal, including
for the linear connections used below. For either a compact smooth
connection or a linear connection, $|a(z)|\le c_a(1+|z|)$ and
$a$ is Lipschitz. Choose
$s(z)=K|z|^{-2}e^{-M|z|/2}$ so that
$|S_M(z)|(1+|z|)^{24}\le s(z)$. Its convolution powers obey
\begin{equation}\label{eq:probe-convolutions}
 s^{*2}(z)\le C|z|^{-1}e^{-M|z|/2},\quad
 s^{*2}(z)\le C|z|^{-3/2},\quad
 s^{*3}(z)\le C|z|^{-1/2}.
\end{equation}
For example, the bound for the second convolution follows by factoring
out $e^{-M|z|/2}$ and scaling
$\int|z-w|^{-2}|w|^{-2}\,\dd w=C|z|^{-1}$; the third follows by
convolving $|z|^{-2}$ and $|z|^{-3/2}$. Thus $s^{*2}\in L^2$,
$s^{*4}=s^{*2}*s^{*2}$ is bounded and continuous, and so is
$s^{*5}=s*s^{*4}$.

Let $R_a^{(j)}(x,y)$ be the integral kernel of
$(-1)^j(S_MV_a)^jS_M$, initially off the diagonal, for $0\le j\le4$.
Distributing the factors $1+|z_l|$ along the chain of propagators gives
\begin{equation}\label{eq:probe-kernel-majorant}
 |R_a^{(j)}(x,y)|\le[c_a(1+|x|)]^j s^{*(j+1)}(x-y).
\end{equation}
These integrals act as the indicated operators on Schwartz spinors;
absolute integration against such a spinor justifies every composition.
The fourth kernel extends continuously to all $(x,y)$. To see this,
use the four independent increments of the first three and last
propagators. Changing $(x,y)$ then changes four Lipschitz factors and
the middle propagator. On $|x|,|y|,|x'|,|y'|\le r$, the resulting
difference is bounded by
\begin{multline}\label{eq:probe-kernel-continuity}
 C_{a,r}\bigl((|x-x'|+|y-y'|)\|s^{*5}\|_\infty\\*
 +\|s^{*4}\|_\infty
       \|S_M-S_M(\,\cdot+(x'-y')-(x-y))\|_{L^1}\bigr).
\end{multline}
This tends to zero. Define
$r_4(x;M,a)=\tr R_a^{(4)}(x,x)$ using this extension.

For compact $a$, differentiation in $M$ now gives
\begin{equation}\label{eq:probe-mass-derivative}
 \frac{\dd}{\dd M}q_{4,M}(a)=\int r_4(x;M,a)\,\dd x,
 \qquad q_{4,M}(a)\longrightarrow0\quad(M\longrightarrow\infty).
\end{equation}
To justify this differentiation, write the trace as the absolutely convergent
four-vertex integral
\[
 -\tfrac14\int\tr\bigl[S_M(z_4-z_1)V_a(z_1)
 S_M(z_1-z_2)V_a(z_2)S_M(z_2-z_3)V_a(z_3)
 S_M(z_3-z_4)V_a(z_4)\bigr]\,\dd z_1\cdots\dd z_4.
\]
Its absolute integral is at most
$4\|a\|_\infty^3\|a\|_{L^1}s^{*4}(0)$ before the factor $1/4$.
On any compact interval of positive masses, $S_M$ and
$\partial_MS_M$ have a common majorant of the same form. Since
$\partial_MS_M=-S_M*S_M$, the four differentiated terms have equal
integrals by cyclicity. The factor $-1/4$ and the minus sign in this
derivative cancel. Inserting the convolution yields the diagonal
five-propagator integral in \eqref{eq:probe-mass-derivative}; its absolute
integral is bounded by $4\|a\|_\infty^3\|a\|_{L^1}s^{*5}(0)$.
This also proves continuous differentiability. The limit at infinity
follows from \eqref{eq:probe-trace-bound}.

\subsection{The local large-mass expansion}

Write
\[
 V_j(x;w)=i\gamma\cdot\sum_{|\alpha|=j}
                  \frac{w^\alpha}{\alpha!}\partial^\alpha a(x).
\]
The substitution $z_l=x+M^{-1}w_l$ in the diagonal kernel gives
\begin{multline}\label{eq:probe-scaled-kernel}
 r_4(x;M,a)=M^{-2}\tr\int S_1(-w_1)V_a(x+M^{-1}w_1)
 S_1(w_1-w_2)V_a(x+M^{-1}w_2)\\
 {}\cdot S_1(w_2-w_3)V_a(x+M^{-1}w_3)
 S_1(w_3-w_4)V_a(x+M^{-1}w_4)S_1(w_4)\,\dd w.
\end{multline}
Let $\rho_j(x;a)$ be the sum of these integrals with the four
connection factors replaced by $V_{j_l}(x;w_l)$ and
$j_1+\cdots+j_4=j$. Taylor expansion through total order five gives
\begin{equation}\label{eq:probe-local-expansion}
 r_4(x;M,a)=\sum_{j=0}^5M^{-2-j}\rho_j(x;a)+O_a(M^{-8}),
 \qquad M\ge1,
\end{equation}
uniformly in $x$. The remainder contains four factors bounded by
$C\|a\|_{C^6}(1+|w_l|)^6$ and at least the factor $M^{-6}$.
With $w_0=0$,
$\prod_l(1+|w_l|)^6\le\prod_l(1+|w_l-w_{l-1}|)^{24}$.
Use the weighted majorant at unit mass for the five factors $S_1$.
Its fifth convolution at zero bounds the remainder integral by a
constant independent of $M$.
The coefficients vanish off $\supp a$. There is also a spatial tail
bound
\begin{equation}\label{eq:probe-spatial-tail}
 \int_{\operatorname{dist}(x,\supp a)\ge1}|r_4(x;M,a)|\,\dd x
 \le C_ae^{-M}.
\end{equation}
Indeed, the first and last propagators each travel at least
$d=\operatorname{dist}(x,\supp a)$; their bound is $Me^{-Md}$ for
$d\ge1$. Integrating the three intervening propagators gives
\[
 |r_4(x;M,a)|\le108\|a\|_\infty^4|\supp a|M^{-1}e^{-2Md}.
\]
For $M,d\ge1$, $2Md\ge M+d$, which proves
\eqref{eq:probe-spatial-tail}.

The reduction to field strength must be justified at the level of this
diagonal kernel. Let $\lambda\in C_c^\infty(\R^3)$ be real. On Schwartz
spinors, as a formal power series in $u$,
\[
 (M+\gamma\cdot\partial+ui\gamma\cdot(a+d\lambda))^{-1}
 =e^{-iu\lambda}(M+\gamma\cdot\partial+ui\gamma\cdot a)^{-1}
                       e^{iu\lambda}.
\]
All coefficients preserve Schwartz space. Comparing the coefficient
of $u^4$ gives, almost everywhere off the diagonal,
\begin{equation}\label{eq:probe-gauge-kernel}
 R_{a+d\lambda}^{(4)}(x,y)=\sum_{j=0}^4R_a^{(j)}(x,y)
       \frac{[-i(\lambda(x)-\lambda(y))]^{4-j}}{(4-j)!}.
\end{equation}
For $j=0,1,2,3$, the bounds in
\eqref{eq:probe-convolutions}--\eqref{eq:probe-kernel-majorant} show
that these terms are respectively $O(|x-y|^2)$,
$O(|x-y|^2)$, $O(|x-y|^{3/2})$, and $O(|x-y|)$ locally in $x$.
Continuity of the fourth kernel therefore implies
\[
 r_4(x;M,a+d\lambda)=r_4(x;M,a).
\]
Uniqueness of the finite asymptotic expansion in
\eqref{eq:probe-local-expansion} gives the same invariance for
$\rho_j$, $0\le j\le5$.

Fix $x$. If $p^{(k)}(z)$ is the homogeneous part of degree $k$ of
the fifth-order Taylor polynomial of $a(x+z)$, choose a compact
$\lambda$ equal near $x$ to
\[
 \lambda(x+z)=-\sum_{k=0}^5\frac{z\cdot p^{(k)}(z)}{k+1}.
\]
Euler's identity gives $z\cdot(a+d\lambda)^{[5]}(x+z)=0$ for
the Taylor polynomial. Thus the transformed constant term vanishes,
and its first derivative is the antisymmetric matrix
$F_{\nu\mu}(x)/2$, where
$F_{\nu\mu}=\partial_\nu a_\mu-\partial_\mu a_\nu$.
Each term defining $\rho_j$ has four connection factors. Hence it
vanishes for $j<4$, and at order four only $(j_1,j_2,j_3,j_4)=(1,1,1,1)$
remains. Put
\[
 A_{F,\mu}(z)=\tfrac12\sum_\nu z_\nu F_{\nu\mu},\qquad
 \Psi(F)=r_4(0;1,A_F).
\]
Then $\rho_4(x;a)=\Psi(F(x))$. The odd coefficients vanish as well:
substitution $w_l\mapsto-w_l$ and conjugation by $\gamma_5$ leave
the trace unchanged and multiply it by $(-1)^j$. In particular
$\rho_5=0$. Consequently
\begin{equation}\label{eq:probe-field-strength-reduction}
 r_4(x;M,a)=M^{-6}\Psi(F(x))+O_a(M^{-8}).
\end{equation}

\subsection{The constant-field coefficient}

Rotation of the variables and conjugation by the spin lift in $SU(2)$
show that $\Psi(F)$ depends on the vector
$B_l=\tfrac12\sum_{\nu,\mu}\varepsilon_{l\nu\mu}F_{\nu\mu}$ only
through $|B|$. Take $B=(0,0,b)$, $b\ge0$, and
$A=(-bz_2/2,bz_1/2,0)$. Set $D_u=\partial+iuA$ and
\[
 h_u=1-\Delta+uW_1+u^2W_2,\qquad
 W_1=-2iA\cdot\partial+b\sigma_3,\quad
 W_2=|A|^2.
\]
The Clifford relations give
\begin{equation}\label{eq:probe-squared-dirac}
 (1-\gamma\cdot D_u)(1+\gamma\cdot D_u)
   =\diag(h_u,h_u),\qquad
 (1+\gamma\cdot D_u)^{-1}
   =(1-\gamma\cdot D_u)\diag(h_u^{-1},h_u^{-1}).
\end{equation}
These are identities of formal series acting on Schwartz spinors;
no expansion of an operator with an unbounded connection is being
assumed in operator norm.

We compute the fourth coefficient of $h_u^{-1}$ from the kernel
\begin{multline}\label{eq:probe-Mehler}
 k_u(t;x,y)=(4\pi t)^{-3/2}e^{-|x-y|^2/(4t)}
 \frac{\tau}{\sinh\tau}\exp\left(-\frac{(\tau\coth\tau-1)|x_p-y_p|^2}{4t}
                           +\frac{iub}{2}x\wedge y\right)
 e^{-\tau\sigma_3},\\
 \tau=ubt,\quad x_p=(x_1,x_2),\quad x\wedge y=x_1y_2-x_2y_1.
\end{multline}
It is analytic in $u$ for $|u|bt<\pi$. The kernel satisfies
$(\partial_t+h_u-1)k_u=0$. The planar verification uses
$\alpha=(ub/4)\coth(ubt)$ and
$\phi=-\alpha|x_p-y_p|^2+iub(x\wedge y)/2$:
\[
 \frac{\check h e^\phi}{e^\phi}
   =4\alpha-(4\alpha^2-u^2b^2/4)|x_p-y_p|^2,
 \qquad \alpha'=-(4\alpha^2-u^2b^2/4),
\]
where $\check h=-\partial_1^2-\partial_2^2
-iub(x_1\partial_2-x_2\partial_1)+u^2b^2|x_p|^2/4$.
Differentiation of the prefactor supplies $-4\alpha$; differentiation
of $e^{-\tau\sigma_3}$ supplies the spin term. The longitudinal
Gaussian satisfies the one-dimensional heat equation.

Here only the coefficients of \eqref{eq:probe-Mehler} are needed.
Write $k_u=\sum_{n\ge0}u^nk^{(n)}$. Each coefficient is the Gaussian
times a finite polynomial in $t,x-y,y$, whose monomials
$t^j(x-y)^\alpha y^\beta$ have $2j+|\alpha|\ge n$.
The factor $e^{iub(x-y)\wedge y/2}$ has weight at least one per
power of $u$, and the other factors have weight at least two,
if $t$ has weight two and $x-y$ weight one.
Thus
\[
 K^{(n)}(x,y)=\int_0^\infty e^{-t}k^{(n)}(t;x,y)\,\dd t
\]
converges off the diagonal and, for $n\ge1$, on the diagonal.
Its terms are $y^\beta$ times
\[
 m_{j,\alpha}(x-y)=\int_0^\infty e^{-t}t^j(x-y)^\alpha
                 (4\pi t)^{-3/2}e^{-|x-y|^2/(4t)}\,\dd t.
\]
Their Fourier transforms are
$(i\partial_p)^\alpha[j!(1+|p|^2)^{-j-1}]$, so they define
operators preserving Schwartz space.

Coefficientwise integration of the heat equation identifies these
operators with the formal inverse of $h_u$. Indeed, for a Schwartz
spinor $\varphi$, the zeroth heat coefficient converges to $\varphi$
as $t\downarrow0$. The $n$th coefficient for $n\ge1$ tends to zero
uniformly, since every polynomial term is bounded after integration
against $\varphi$ by $C_\varphi t^{j+|\alpha|/2}$.
Integrating from $\epsilon$ to infinity and sending $\epsilon$ to
zero therefore gives
\[
 (1-\Delta)K^{(n)}+W_1K^{(n-1)}+W_2K^{(n-2)}=\delta_{n0}1.
\]
The convergence can be taken in tempered distributions; the
polynomial differential operators are continuous there. Since the
constant operator $1-\Delta$ is invertible on Schwartz space, the
right inverse is the unique formal inverse.

To evaluate the diagonal, the kernel of its fourth coefficient must
agree with the continuous kernel already constructed. The functions
$m_{j,\alpha}$ are continuous when $2j+|\alpha|\ge2$ and continuously
differentiable when $2j+|\alpha|\ge3$. For this, write
$m_{j,\alpha}(z)=z^\alpha\mu_j(|z|)$ with
\[
 \mu_j(s)=\int_0^\infty e^{-t}t^j(4\pi t)^{-3/2}e^{-s^2/(4t)}\,\dd t.
\]
For $j\ge1$, $\mu_j$ is bounded and continuous at zero and
$\mu_j'(s)=-s\mu_{j-1}(s)/2$ for $s>0$; for $j=0$ use
$\mu_0(s)=e^{-s}/(4\pi s)$. These formulas give the asserted
regularity at zero. Consequently $K^{(3)},K^{(4)}$ and
$\partial_xK^{(4)}$ are continuous. From
\eqref{eq:probe-squared-dirac} the fourth Dirac coefficient has kernel
\[
 (1-\gamma\cdot\partial_x)\diag(K^{(4)},K^{(4)})(x,y)
       -i\gamma\cdot A(x)\diag(K^{(3)},K^{(3)})(x,y).
\]
It equals $R_A^{(4)}$ because both continuous kernels give the same
operator on Schwartz space. At $x=y=0$, $A(0)=0$, and the derivative
term has zero trace: the two Pauli blocks have opposite signs.
It follows that
\begin{align}
 \Psi(F)&=2\int_0^\infty e^{-t}\tr k^{(4)}(t;0,0)\,\dd t\\
 &=-\frac4{45}(4\pi)^{-3/2}\Gamma(7/2)b^4
   =-\frac1{48\pi}|B|^4.\label{eq:probe-constant-field}
\end{align}
In the second line we used
$\tr k_u(t;0,0)=2(4\pi t)^{-3/2}\tau\coth\tau$ and
$\tau\coth\tau=1+\tau^2/3-\tau^4/45+O(\tau^6)$.

\subsection{Integration and dilation}

Integrate \eqref{eq:probe-field-strength-reduction} over the unit
neighborhood of $\supp a$, and use \eqref{eq:probe-spatial-tail}
on its complement. Equations \eqref{eq:probe-mass-derivative} and
\eqref{eq:probe-constant-field} give
\[
 \frac{\dd}{\dd M}q_{4,M}(a)
 =-\frac1{48\pi M^6}\int|\operatorname{curl}a|^4\,\dd x
                                    +O_a(M^{-8}).
\]
Integration from $M$ to infinity proves
\eqref{eq:large-mass-quartic}.

Finally let $\varphi(x)=e^{-1/(1-|x|^2)}$ for $|x|<1$ and zero
otherwise, set $\beta_{12}=\varphi$, $\beta_{13}=\beta_{23}=0$, and
$a=d^*\beta=(\partial_2\varphi,-\partial_1\varphi,0)$.
The third component of $\operatorname{curl}a$ is
$-(\partial_1^2+\partial_2^2)\varphi$ and has value $4e^{-1}$ at zero.
Thus $I=\int|\operatorname{curl}a|^4\,\dd x>0$.
For $\ell\ge1$ put
$\beta^{(\ell)}(x)=\ell^2\beta(x/\ell)$,
$a_\ell(x)=\ell a(x/\ell)$, and $f_\ell=dd^*\beta^{(\ell)}$.
The unitary dilation $U_\ell\psi(x)=\ell^{-3/2}\psi(x/\ell)$ gives
\[
 U_\ell^{-1}S_{\bar m}V_{a_\ell}U_\ell
       =\ell^2S_{\bar m\ell}V_a,
 \qquad q_{4,\bar m}(a_\ell)=\ell^8q_{4,\bar m\ell}(a).
\]
For $\bar m\ell\ge1$, therefore,
\begin{equation}\label{eq:probe-positive-dilation}
 b(f_\ell)=24q_{4,\bar m}(a_\ell)
   =\frac{\ell^3I}{10\pi\bar m^5}
                      +O_a(\ell\bar m^{-7}).
\end{equation}
This is positive for sufficiently large $\ell$ and completes the proof
of \cref{prop:positive-probe}.


\section{The chamber symbol and determinant derivatives}\label{subsec:response-bloch}\label{sec:free-response}

The free chamber symbol has a single four-component Dirac singularity.
We first isolate it, then use its inverse bounds to control the contact
terms in the fourth derivative.

\subsection{The singular part of the inverse}

Use the fiber $V=\ell^2(O_{\mathrm{cub}};\C^4)$ with measure
$48^{-1}\sum_g$ and the constant-label isometry $(Jz)(g)=z$.
The orthogonal projection onto its complement is $Q=1-JJ^*$.
The fiber $V$ has dimension $192$, and traces use this full dimension;
the label measure does not multiply the matrix trace by
$1/48$. Write $\mathcal D_m(p)$ for the free symbol of
\cite[\S2.2]{FermionBlock}, and $R_m(p)=\mathcal D_m(p)^{-1}$.
For real $p$ it is invertible when $m>0$.

The light and heavy blocks are the blocks relative to
$\operatorname{ran}J\oplus\operatorname{ran}Q$. The heavy inverse,
the vanishing of the mixed blocks at zero momentum, and the
ellipticity of the light Schur complement are proved in
\cite[Proposition~2.2]{FermionBlock}.

We need a bound on the difference between the full inverse
and its single light Dirac part, including mixed and heavy
blocks. There is $C_\rho<\infty$ such that
\begin{equation}\label{eq:response-bloch-remainder}
 R_m(p)=J(m+i\gamma\cdot p)^{-1}J^*+\mathcal E_m(p),\qquad
 \|\mathcal E_m(p)\|\le C_\rho,\qquad p\in[-\pi,\pi]^3,\quad0<m\le1.
\end{equation}
To prove this, use the block inversion formula
\cite[(16)]{FermionBlock}. The off-diagonal blocks
vanish at zero momentum, and are $O(|p|)$ because only
the external hop varies with $p$. The heavy inverse is
bounded by \cite[(10)]{FermionBlock}. The light Schur complement
therefore has the form
\[
 B_m(p)=m+i\gamma\cdot p+r_m(p),\qquad
                    \|r_m(p)\|\le C_\rho|p|^2.
\]
For $|p|$ below a fixed $\rho$-dependent radius, the product
$r_m(p)(m+i\gamma\cdot p)^{-1}$ has norm at most one half.
Its Neumann inverse gives
\[
 \|B_m(p)^{-1}-(m+i\gamma\cdot p)^{-1}\|
        \le2C_\rho |p|^2(m^2+|p|^2)^{-1}\le2C_\rho,
 \qquad |p|\|B_m(p)^{-1}\|\le2.
\]
In the other four terms of the block inverse, each mixed
block supplies a factor $|p|$ beside this light inverse.
The heavy terms are bounded. This proves
\eqref{eq:response-bloch-remainder} near zero. Away from
zero, the real-part and high-momentum ellipticity bounds
of \cite[\S2.2]{FermionBlock} bound the full inverse uniformly,
and the continuum inverse is bounded by $1/|p|$. This
proves the estimate on the whole cube. In particular
\begin{equation}\label{eq:response-bloch-bounds}
 \|R_m(p)\|\le C_\rho+(m^2+|p|^2)^{-1/2}
                    \le\frac{C'_\rho}{m+|p|}.
\end{equation}

Define external-hop matrices
\[
 (T_\nu^\pm u)(g)=1_{\{ge_1=\pm e_\nu\}}
                          3(\rho\mp\gamma_\nu)u(gr_{e_1}).
\]
Then
\begin{align}\label{eq:response-symbol-vertices}
 \mathcal D_m(p)&=\mathcal D_{m,\mathrm{loc}}
       -\sum_\nu(e^{ip_\nu}T_\nu^++e^{-ip_\nu}T_\nu^-),\notag\\
 V_\nu(p):=\partial_\nu\mathcal D_m(p)
       &=-i(e^{ip_\nu}T_\nu^+-e^{-ip_\nu}T_\nu^-),\notag\\
 W_\nu(p):=\partial_\nu^2\mathcal D_m(p)
       &=e^{ip_\nu}T_\nu^++e^{-ip_\nu}T_\nu^-.
\end{align}
Mixed second derivatives vanish. The disjoint rows indexed
by $ge_1=\pm e_\nu$ imply
$\|\sum z_\nu^\pm T_\nu^\pm\|\le3(1+\rho)\max|z_\nu^\pm|$.
Since eight labels contribute to each direction,
\begin{equation}\label{eq:response-light-vertex}
 J^*T_\nu^\pm J=\tfrac12(\rho\mp\gamma_\nu),\qquad
 J^*V_\nu(0)J=i\gamma_\nu.
\end{equation}
All nonzero derivatives of the symbol are bounded by six
on real momenta. Differentiating $R_m\mathcal D_m=1$ gives
\[
 \partial^\alpha R_m=-\sum_{0<\delta\le\alpha}
   \binom\alpha\delta(\partial^{\alpha-\delta}R_m)
                    (\partial^\delta\mathcal D_m)R_m.
\]
Starting with \eqref{eq:response-bloch-bounds}, induction bounds
each summand by a constant times
$(\operatorname{dist}(p,2\pi\Z^3)+m)^{-1-|\alpha|}$;
the remaining factor of the distance has nonnegative exponent
and is bounded by a power of seven. Thus
\begin{equation}\label{eq:response-inverse-derivatives}
 \|\partial^\alpha R_m(p)\|
   \le C_{\alpha,\rho}
       (\operatorname{dist}(p,2\pi\Z^3)+m)^{-1-|\alpha|}.
\end{equation}
The involution $(\mathcal Uu)(g)=\gamma_5u(-g)$ satisfies
\begin{equation}\label{eq:response-bloch-parity}
 \mathcal U\mathcal D_m(p)\mathcal U=\mathcal D_m(-p),\qquad
 \mathcal UV_\nu(p)\mathcal U=-V_\nu(-p).
\end{equation}
Indeed $-g$ reverses every bond direction, and
$\gamma_5\gamma_\nu\gamma_5=-\gamma_\nu$ restores the
hopping matrix. The external phase then changes $p$ to
$-p$. This parity will cancel the first momentum derivative
of the polarization kernel.

\subsection{Derivative and contact terms}\label{subsec:response-contact}

Let $\mathcal P_n=\frac{2\pi}{n}(\Z+\frac12)^3\cap(-\pi,\pi]^3$.
The normalized Bloch vectors are $n^{-3/2}e^{ip\cdot c}$,
and the lattice free propagator $G$ is the multiplier
$R_m(p)$ on these antiperiodic momenta. Write
$A_\nu(c)=a_{\eta,\nu}(\eta c)$ and
$A_\nu^-(c)=A_\nu(c-e_\nu)$, and let
$U_\nu^\pm\psi(c)=\psi(c\pm e_\nu)$. The bond values are
$b_\nu^+=\eta A_\nu$, $b_\nu^-=-\eta A_\nu^-$.
The exact connection derivatives are
\begin{equation}\label{eq:response-connection-derivatives}
 D_j=-\sum_{\nu,\pm}(ib_\nu^\pm)^jT_\nu^\pm U_\nu^\pm,
 \quad D^{\mathrm{ch}}[B_t]-D^{\mathrm{ch}}[0]
                  =\sum_{j\ge1}\frac{t^j}{j!}D_j,
 \quad Y_j=GD_j.
\end{equation}
The shifts act first. These $D_j$ are derivatives, with
the factorials outside the vertices.

For small complex $t$, $N_0(t)=\det(1+\sum_{j\ge1}t^jY_j/j!)$
and its zero-at-zero logarithm is the convergent finite-matrix
trace series $\sum_{l\ge1}(-1)^{l+1}\operatorname{Tr}(\sum_j
t^jY_j/j!)^l/l$. The tail $l\ge5$ is $O(t^5)$, so it
contributes no fourth derivative. Cyclicity of the trace
in the terms $l\le4$ gives
\begin{equation}\label{eq:response-fourth-identity}
 \begin{split}
 b_{\mathrm{ch}}={}&\operatorname{Tr}Y_4
       -4\operatorname{Tr}(Y_1Y_3)-3\operatorname{Tr}Y_2^2\\
       &+12\operatorname{Tr}(Y_1^2Y_2)-6\operatorname{Tr}Y_1^4.
 \end{split}
\end{equation}
For example, the cyclically equal orders $(1,1,2)$,
$(1,2,1)$ and $(2,1,1)$ contribute
$24\cdot\frac13\cdot\frac32=12$. The quadratic identity
is $a_{\mathrm{ch}}=\operatorname{Tr}Y_2-\operatorname{Tr}Y_1^2$.
The determinant series is used only near zero; the determinant
itself is entire in $t$.

We bound all four contact traces in
\eqref{eq:response-fourth-identity} uniformly in $R$ and the
position of the probe. The interpolating periodic field is
\[
 \widetilde A_\nu(y)=\sum_z
             a_{\eta,\nu}^0(y+\eta e_\nu/2+Rz),
 \qquad A_\nu(c)=\widetilde A_\nu(\eta c).
\]
When $R\ge D_\beta+1$, averaging translates of the derivatives
of $\beta$ shows
\begin{equation}\label{eq:response-probe-estimates}
 \begin{gathered}
 \|\widetilde A\|_\infty+\|\partial\widetilde A\|_\infty\le C_\beta^{\mathrm{free}},
 \qquad \|\widetilde A-a^R\|_\infty\le C_\beta^{\mathrm{free}}\eta,\\
 \eta^3|\{c:A_\nu(c)\ne0\text{ for some }\nu\}|\le(D_\beta+3)^3,
 \qquad
 \|\eta A_\nu\|_{\ell^4}\le C_\beta^{\mathrm{free}}\eta^{1/4},\\
 \|\eta(A_\nu^--A_\nu)\|_{\ell^4}\le C_\beta^{\mathrm{free}}\eta^{5/4},
 \qquad\|\widetilde A-a^R\|_{L^4(T_R)}\le C_\beta^{\mathrm{free}}\eta.
 \end{gathered}
\end{equation}
Every averaging displacement has size at most $\eta$ in
sup norm. The translates of $\supp\beta$ are disjoint,
so their derivative sup norms are unchanged. The supports
of the interpolants lie in $Q_\beta+[-\eta,\eta]^3$ modulo
the periods. Its volume and number of lattice points give
the last estimates, including the union with its one-step
translate for the difference.

Let $\widehat{A_\nu^j}(k)=n^{-3}\sum_cA_\nu(c)^je^{-ik\cdot c}$
on the finite momentum group $\mathcal G_n=\frac{2\pi}{n}\Z^3/
2\pi\Z^3$, and set $F_j(k)=\sum_\nu|\widehat{A_\nu^j}(k)|$.
There are constants $C_{\mathrm{Four},1},C_{\mathrm{Four},0}$ depending only on $\beta$ with
\begin{equation}\label{eq:response-probe-fourier-norms}
 \|F_j\|_{\ell^1(\mathcal G_n)}\le C_{\mathrm{Four},1}^j,
 \qquad R^3\|F_j\|_\infty\le C_{\mathrm{Four},1}^{j-1}C_{\mathrm{Four},0}.
\end{equation}
The transform of the centered difference and the cube average
in \eqref{eq:response-lattice-probe} is
\[
 -\sum_{\mu\ne\nu}iP_\mu\operatorname{sinc}(\eta P_\mu/2)
 \prod_l\operatorname{sinc}(\eta P_l/2)\widehat\beta_{\mu\nu}(P),
\]
where $\operatorname{sinc}(z)=\sin z/z$ and $\operatorname{sinc}(0)=1$.
The midpoint shift multiplies this transform by a phase of modulus one.
The shifted transform has modulus at most
$C_\beta^{\mathrm{free}}(1+|P|)^{-4}$. Poisson summation \cite[Theorems~4--7]{GuthFourier2015} gives
\[
 \widehat A_\nu(k)=R^{-3}\sum_{z\in\Z^3}
      \widehat{a_{\eta,\nu}^0(\,\cdot+\eta e_\nu/2)}
                                      ((k+2\pi z)/\eta).
\]
Summing over representatives $k$ and $z$ gives a single
sum over $\frac{2\pi}{R}\Z^3$, bounded after multiplication
by $R^{-3}$ by its integrable radial majorant, uniformly
for $R\ge1$. For $|k|_\infty\le\pi$, nonzero $z$ have
$|(k+2\pi z)/\eta|\ge\pi|z|_\infty/\eta$; the convergent
sum $\sum_{z\ne0}|z|_\infty^{-4}$ gives the sup bound.
Powers of $A_\nu$ give convolutions on $\mathcal G_n$;
$\ell^1*\ell^1\to\ell^1$ and $\ell^1*\ell^\infty\to\ell^\infty$
give \eqref{eq:response-probe-fourier-norms} for all $j$.

The momentum blocks of a vertex are
\[
 D_j(p,p')=-\eta^j\sum_\nu\widehat{A_\nu^j}(p-p')
       [i^je^{ip'_\nu}T_\nu^++(-i)^je^{-ip_\nu}T_\nu^-],
 \qquad\|D_j(p,p')\|\le6\eta^jF_j(p-p').
\]
For $l\le3$ and $j_1+\cdots+j_l=4$, apply these bounds
and \eqref{eq:response-bloch-bounds} to the cyclic trace.
With differences $k_i=p_i-p_{i+1}$, the sum is parametrized
by $p_1,k_1,\ldots,k_{l-1}$. H\"older in $p_1$ bounds
the product of its $l$ translates of $(m+|p|)^{-1}$ by
$\sum_p(m+|p|)^{-l}$. Bound the final $F_{j_l}$ in sup
norm and the others in $\ell^1$. Thus
\begin{equation}\label{eq:response-cyclic-contact}
 |\operatorname{Tr}(Y_{j_1}\cdots Y_{j_l})|
   \le192(6C'_\rho)^l C_{\mathrm{Four},1}^3C_{\mathrm{Four},0}\,
          \eta\left(n^{-3}\sum_{p\in\mathcal P_n}(m+|p|)^{-l}\right).
\end{equation}
The factor $\eta$ is $\eta^4n^3/R^3$.

For antiperiodic momenta, comparison with their disjoint
cubes of side $2\pi/n$ gives, for a decreasing radial
function $h$, $n^{-3}\sum_ph(|p|)\le\pi^{-3}\int h(|y|)\,\dd y$.
Indeed in the cube of $p$ one has $|y|\le2|p|$, and a
rescaling by two gives the factor $\pi^{-3}$. On the
bounded momentum cube this yields constants for $l=1,2$
and $\log(1+5/m)$ for $l=3$. Inserting the multiplicities
$1,4,3,12$ from \eqref{eq:response-fourth-identity} proves
\begin{equation}\label{eq:response-contact-bound}
 \begin{split}
 &|\operatorname{Tr}Y_4|+4|\operatorname{Tr}(Y_1Y_3)|
          +3|\operatorname{Tr}Y_2^2|+12|\operatorname{Tr}(Y_1^2Y_2)|\\
 &\hspace{25mm}\le C_{\mathrm{ct}}(\rho,\beta)\,
               \eta\left(1+\log(1+5/(\bar m\eta))\right).
 \end{split}
\end{equation}
The constant does not depend on $\bar m$, $\eta$, $R$ or
the position of $\supp\beta$.

\section{Comparison of the fourth traces}\label{subsec:response-fourth}

The contact bound reduces the first estimate in
\eqref{eq:response-fourth-errors} to the trace with four first-order
vertices. We compare that trace with the continuum torus trace in
five steps. The Schatten-four norm controls each replacement,
and the Yukawa kernel controls the remaining volume error.

\begin{proof}[Proof of the fourth-coefficient assertions in \cref{prop:response-fourth}]
We first establish the trace bounds used in the comparisons.
For a lattice multiplier $M_h$ with fiber symbol $h(p)$
and a scalar multiplication operator $v$, one has
\begin{equation}\label{eq:response-four-ideal-lattice}
 \|M_hv\|_4\le
       \left(n^{-3}\sum_{p\in\mathcal P_n}\|h(p)\|_4^4\right)^{1/4}
                                         \|v\|_{\ell^4}.
\end{equation}
To prove it put $U=M_{h^*h}$ and $V=|v|^2$. Expanding the
nonnegative trace of $(UV-VU)^*(UV-VU)$ gives
$\operatorname{Tr}(UVUV)\le\operatorname{Tr}(U^2V^2)$.
The left side is $\|M_hv\|_4^4$, and each normalized
Bloch vector has constant squared modulus $n^{-3}$.
Evaluating the right side in that basis gives
\eqref{eq:response-four-ideal-lattice}. On the continuum
torus the same proof, first for a finite set of momenta,
gives
\begin{equation}\label{eq:response-four-ideal-torus}
 \|M_hv\|_4\le
       \left(R^{-3}\sum_{P\in\mathcal P_R}\|h(P)\|_4^4\right)^{1/4}
                                           \|v\|_{L^4(T_R)}.
\end{equation}
Truncation in momentum converges in the Schatten norm
when the sum is finite, and its operator-norm limit is
$M_hv$, proving the general form.

The cube comparison used above also gives
\begin{equation}\label{eq:response-four-momentum-sums}
 \begin{gathered}
 n^{-3}\sum_{p\in\mathcal P_n}(m^2+|p|^2)^{-2}\le\frac1{\pi m},
 \qquad
 R^{-3}\sum_{P\in\mathcal P_R}(\bar m^2+|P|^2)^{-2}
                                             \le\frac1{\pi\bar m},\\
 R^{-3}\sum_{|P|_\infty\le\pi/(2\eta)}|P|^{-2}\le C\eta^{-1},
 \qquad
 R^{-3}\sum_{|P|_\infty>\pi/(2\eta)}|P|^{-4}\le C\eta.
 \end{gathered}
\end{equation}
For the last bound use the decreasing radial majorant
$\min((2\eta/\pi)^4,|P|^{-4})$; its integral is $O(\eta)$.
The preceding bounds and
$\|s_{\bar m}(P)\|_4^4=4(\bar m^2+|P|^2)^{-2}$ show that
$S^Ri\gamma\cdot a^R$ has a bounded Schatten-four norm,
uniformly in $R\ge D_\beta+1$. The same estimate applies
to the interpolating field $\widetilde A$.

On $\R^3$, let
\[
 Y_{\bar m}(x)=\frac{e^{-\bar m|x|}}{4\pi|x|},\qquad
 S_{\bar m}(x)=(\bar m-\gamma\cdot\partial)Y_{\bar m}(x).
\]
The scalar Fourier transform is $(\bar m^2+|P|^2)^{-1}$;
the derivative has no point mass, since its boundary
term on a sphere of radius $r$ is $O(r)$. Its Fourier
transform therefore gives $s_{\bar m}(P)$. These kernels
are integrable, so their convolutions are the corresponding
bounded Fourier multipliers on $L^2$. Also
$\|Y_{\bar m}\|_{L^2}^2=(8\pi\bar m)^{-1}$.
For bounded compactly supported $v$, the kernel of
$v^*S_{\bar m}^*S_{\bar m}v$ is
$\overline{v(x)}Y_{\bar m}(x-y)v(y)1_4$. Tonelli and
Cauchy--Schwarz in the translated factors give
\[
 \|S_{\bar m}v\|_4^4
   \le4\|Y_{\bar m}\|_{L^2}^2\|v\|_{L^4}^4
   =\frac{\|v\|_{L^4}^4}{2\pi\bar m}.
\]
For the matrix multiplication $i\gamma\cdot a$, the
Clifford identity gives the same proof with $|a|$ in
place of $|v|$. Thus
\begin{equation}\label{eq:response-four-ideal-infinite}
 \|S_{\bar m}i\gamma\cdot a\|_4^4\le
                    (2\pi\bar m)^{-1}\int_{\R^3}|a|^4\,\dd x.
\end{equation}
The traces in
\eqref{eq:response-continuum-fourth} exist by Schatten
H\"older.

To check reality, treat $S=S_{\bar m}$ or $S^R$ and
$v=i\gamma\cdot a$ or $i\gamma\cdot a^R$ together.
Factor $v=|a|^{1/2}(i\gamma\cdot a/|a|)|a|^{1/2}$,
with the middle factor zero at its zeros. The preceding
estimates put $SvS$ in the Hilbert--Schmidt class.
One has $\gamma_5S\gamma_5=S^*$ and
$\gamma_5v\gamma_5=v^*$. Hence conjugating
$\operatorname{Tr}[(SvS)v(SvS)v]$ by $\gamma_5$ gives
its complex conjugate after cyclically moving the
bounded $v$ past a trace-class product of the two
Hilbert--Schmidt factors. This proves reality.

We now compare the lattice fourth-power trace with its
continuum torus value. The four-power inequality used
throughout is
\begin{equation}\label{eq:response-fourth-trace-difference}
 |\operatorname{Tr}X^4-\operatorname{Tr}Z^4|
 \le4\max(\|X\|_4,\|Z\|_4)^3\|X-Z\|_4.
\end{equation}
It follows by expanding $X^4-Z^4$ into four products
with one $X-Z$ factor and applying Schatten H\"older.
All the operators compared below have uniformly bounded
Schatten-four norms; each bound follows from
\eqref{eq:response-four-ideal-lattice} or
\eqref{eq:response-four-ideal-torus}. In particular
\eqref{eq:response-bloch-bounds} gives
$n^{-3}\sum_p\|R_m(p)\|_4^4\le
C_\rho(1+m^{-1})$, which, with
$\|\eta A_\nu\|_{\ell^4}\le C_\beta^{\mathrm{free}}\eta^{1/4}$,
proves $\|Y_1\|_4\le C_{\rho,\bar m,\beta}$.

\emph{Removing the bounded part of the inverse.}
Let $\chi(p)=1_{\{|p|_\infty\le\pi/2\}}$ and
$S_c$ be the lattice multiplier $\chi(p)s_m(p)$ on
four-component antiperiodic spinors. The complementary
inverse
$R_m(p)-J\chi(p)s_m(p)J^*$ is uniformly bounded:
\eqref{eq:response-bloch-remainder} controls the whole
cube, and $s_m$ is bounded on $1-\chi$. Thus
\[
 \|Y_1-JS_cJ^*D_1\|_4\le C_{\rho,\beta}\eta^{1/4}.
\]
The two norms in \eqref{eq:response-fourth-trace-difference}
are bounded by the same estimates. Since $J^*J=1$,
cyclicity reduces the fourth-power trace of the light
operator to that of $S_cw_1$, where
\begin{equation}\label{eq:response-compressed-vertex}
 w_1=J^*D_1J=-\frac{i\eta}{2}\sum_\nu
       [A_\nu(\rho-\gamma_\nu)U_\nu^+
             -A_\nu^-(\rho+\gamma_\nu)U_\nu^-].
\end{equation}
This comparison costs $C\eta^{1/4}$ in the trace. The
compression uses \eqref{eq:response-light-vertex}, and
therefore has one four-component light species.

\emph{Replacing the vertex by multiplication.}
Move shifts to the left by
$A_\nu U_\nu^+=U_\nu^+A_\nu^-$ and
$A_\nu^-U_\nu^-=U_\nu^-A_\nu$. Then
\begin{equation}\label{eq:response-vertex-error}
 \begin{split}
 w_1-i\eta\gamma\cdot A=-\frac{i\eta}{2}\sum_\nu\big[&
  (\rho-\gamma_\nu)(U_\nu^+-1)A_\nu^-
       -(\rho+\gamma_\nu)(U_\nu^--1)A_\nu\\
       &+(\rho-\gamma_\nu)(A_\nu^--A_\nu)\big].
 \end{split}
\end{equation}
The multiplier
$\chi(p)s_m(p)(\rho\mp\gamma_\nu)(e^{\pm ip_\nu}-1)$
has norm at most two, since $|e^{ip_\nu}-1|\le|p_\nu|$.
Its Schatten-four multiplier average is consequently
bounded independently of the mass. The first two terms
of \eqref{eq:response-vertex-error} then cost
$O(\eta^{1/4})$. The last term costs
$O(m^{-1/4}\eta^{5/4})=O(\eta)$ by
\eqref{eq:response-probe-estimates}. Therefore
$\|S_c(w_1-i\eta\gamma\cdot A)\|_4\le C\eta^{1/4}$,
and \eqref{eq:response-fourth-trace-difference} gives
the same trace error.

\emph{Identifying the truncated continuum momenta.}
The map $p\mapsto P=p/\eta$ identifies the retained
lattice momenta with
$\mathcal P_R^<=\{P\in\mathcal P_R:|P|_\infty\le\pi/(2\eta)\}$,
and $\eta s_m(p)=s_{\bar m}(P)$. The operator
$S_ci\eta\gamma\cdot A$ has range in this truncated
space. Its fourth trace is thus that of its compression
to the space. In the normalized momentum basis this
compression has blocks
\[
 T_{\mathrm{lat}}(P,P')=s_{\bar m}(P)i\sum_\nu\gamma_\nu
                 n^{-3}\sum_cA_\nu(c)e^{-i\eta(P-P')\cdot c}.
\]
Let $\widetilde A$ be the smooth interpolant and
$\widehat A_R(k)=R^{-3}\int_{T_R}\widetilde A(y)e^{-ik\cdot y}\,\dd y$.
Absolute Fourier convergence, using four derivatives
of the interpolant, gives the exact alias identity
\begin{equation}\label{eq:response-aliasing}
 n^{-3}\sum_cA_\nu(c)e^{-i\eta k\cdot c}
    =\widehat A_{R,\nu}(k)+
          \sum_{z\ne0}\widehat A_{R,\nu}(k+2\pi z/\eta).
\end{equation}
For $k=P-P'$, $|k|_\infty\le\pi/\eta$. Four integrations
by parts give
$|\widehat A_{R,\nu}(k')|\le C_\beta^{\mathrm{free}} R^{-3}|k'|_\infty^{-4}$;
the nonzero aliases in \eqref{eq:response-aliasing} are
therefore bounded by $C_\beta^{\mathrm{free}} R^{-3}\eta^4$.

Replacing the discrete matrix elements by $\widehat A_R$
gives the compression $T_{\mathrm{int}}$ of
$\chi S^Ri\gamma\cdot\widetilde A$. Its norm is bounded.
For the difference, sum the squared Hilbert--Schmidt
norms of its blocks. There are at most $R^3\eta^{-3}$
choices of $P'$, while
$\sum_{P\in\mathcal P_R^<}|P|^{-2}\le CR^3\eta^{-1}$.
Since $\|\gamma\cdot z\|_2^2=4|z|^2$, the alias bound
gives
\[
 \|T_{\mathrm{lat}}-T_{\mathrm{int}}\|_2^2
     \le C_\beta^{\mathrm{free}}(R^3\eta^{-1})(R^3\eta^{-3})
                                     (R^{-3}\eta^4)^2
     =C_\beta^{\mathrm{free}}\eta^4.
\]
All powers of $R$ cancel. The Schatten-four norm is
at most the Hilbert--Schmidt norm, so the trace error
in this comparison is $O(\eta^2)$.

\emph{Replacing the interpolant and removing the cutoff.}
The difference between the compressions of
$\chi S^Ri\gamma\cdot\widetilde A$ and
$\chi S^Ri\gamma\cdot a^R$ has Schatten-four norm
$O(\eta)$ by \eqref{eq:response-four-ideal-torus} and
\eqref{eq:response-probe-estimates}. For an operator
$X$ with $X=\chi X$, cyclicity gives
$\operatorname{Tr}X^4=\operatorname{Tr}(\chi X\chi)^4$;
the same holds for $X=X\chi$. Thus we can remove the
compression and compare $\chi S^Ri\gamma\cdot a^R$
directly with $S^Ri\gamma\cdot a^R$. The removed momentum
tail in \eqref{eq:response-four-momentum-sums} gives
\[
 \|(1-\chi)S^Ri\gamma\cdot a^R\|_4\le C_\beta^{\mathrm{free}}\eta^{1/4}.
\]
These are respectively the fourth and fifth comparisons;
their trace errors are $O(\eta)$ and $O(\eta^{1/4})$.
The preceding three comparisons, together with
\eqref{eq:response-contact-bound}, prove the first line
of \eqref{eq:response-fourth-errors}. Here
$\eta(1+\log(1+5/(\bar m\eta)))\le C_{\bar m}\eta^{1/4}$.

It remains to compare torus and infinite-space traces.
Because the multiplication by $a$ is supported in
$Q_\beta$, both traces can be computed on $L^2(Q_\beta;\C^4)$,
using extension by zero to $\R^3$ or antiperiodic extension
to $T_R$. The same compression identity proves this
localization. Convolution by the integrable kernel
$S_{\bar m}(x)$ on antiperiodic functions agrees with
$S^R$ on its momentum basis and hence on $L^2$.
After restricting its two variables to $Q_\beta$, its
difference from the infinite-space kernel is
\[
 W_R(x,y)=\sum_{z\in\Z^3\setminus\{0\}}
                  (-1)^{z_1+z_2+z_3}S_{\bar m}(x-y-Rz).
\]
For $|u|\ge1$ the Yukawa formula gives
$\|S_{\bar m}(u)\|\le(3/(4\pi))e^{-\bar m|u|}$.
Moreover $|x-y-Rz|\ge R|z|_\infty-D_\beta\ge
(R-D_\beta)+|z|_\infty-1\ge1$. Counting
$24j^2+2$ images at $|z|_\infty=j$ yields
\[
 \|W_R(x,y)\|\le C_{\bar m}e^{-\bar m(R-D_\beta)},
 \qquad x,y\in Q_\beta.
\]
The series converges absolutely and uniformly there.
Multiplication by $i\gamma\cdot a(y)$ then gives a
Hilbert--Schmidt difference with this same exponential
bound, since $Q_\beta$ has finite volume. The infinite
localized operator has bounded Schatten-four norm,
and the torus one is bounded by it and the difference.
Applying \eqref{eq:response-fourth-trace-difference}
proves the second line of \eqref{eq:response-fourth-errors}.

The constants above involve only $\rho$, $\bar m$,
the side of the support cube, derivative sup and
$L^1$ norms of $\beta$, and moduli of its Fourier
transform. All are unchanged by a real translation.
This gives the asserted uniformity and the limit.
\end{proof}

\section{Symmetry and the second-order response}\label{subsec:response-symmetry}

\subsection{Evenness and placement}

The coefficients in \eqref{eq:response-free-ratio} are real, and
\begin{equation}\label{eq:response-coefficient-relations}
 c_{2j+1}=0,\qquad c_2=a_{\mathrm{ch}},\qquad
                 c_4=b_{\mathrm{ch}}+3c_2^2.
\end{equation}
We prove the determinant identities for complex external connections
to obtain evenness for every complex value of $t$.

For a real connection, let
$(R_i(B)\psi)(c,g)=\psi(c,gs_i)$ for $i=0,1,2$, and
\[
 (R_3(B)\psi)(c,g)=e^{iB(c,c+ge_1)}
                         \psi(c+ge_1,gr_{e_1}).
\]
The outgoing unit bond direction is $\widehat d_i(c,g)=-gn_i$.
Each $R_i(B)$ is self-adjoint and unitary. The direction of a bond changes
sign at its other endpoint. Consequently
\[
 D^{\mathrm{ch}}[B]=\frac{\rho}{2}
       \sum_iw_i(1-R_i(B))+\frac12
       \sum_iw_i(\gamma\cdot\widehat d_i)R_i(B),
\]
where the first sum is nonnegative and the second is anti-self-adjoint.
The two spin sectors $\gamma_\mu=\diag(\sigma_\mu,-\sigma_\mu)$
give adjoint restrictions, as in \cite[(17)]{ChamberLattice}. Thus
$\det(D^{\mathrm{ch}}[B]+m)$ is a squared absolute value.
The real part of $D^{\mathrm{ch}}[B]+m$ is at least $m>0$,
so the determinant is strictly
positive. In particular $N_0(t)>0$ for real $t$.

The transpose identity extends charge symmetry to complex connections.
Let $C=\sigma_2\oplus\sigma_2$ be the charge-conjugation matrix of
\cite[\S3.1]{ChamberLattice}. Transposing a phased bond reverses its
direction and changes its phase from $B$ to $-B$.
The identity $C\gamma_v^\top C^{-1}=-\gamma_v$ changes the transposed
matrix at the other endpoint back to the matrix at the original one.
The diagonal part is scalar; hence
\[
 C(D^{\mathrm{ch}}[B]+m)^\top C^{-1}
                   =D^{\mathrm{ch}}[-B]+m.
\]
The transpose here is for the bilinear chamber pairing. In the site
basis this pairing is $1/48$ times the identity, so this is the
ordinary matrix transpose and preserves the determinant. Every matrix
entry is a finite sum of constants and exponentials of linear
functions of $t$. Therefore $N_0$ is entire, and the transpose
identity gives $N_0(-t)=N_0(t)$ for all $t\in\C$.
Expanding its logarithm at zero proves
\eqref{eq:response-coefficient-relations}.

We also record exact covariance under lattice translations and cubic
reflections. Let $g=(v,r)$, where $v\in\Z^3$ and $r$ is a signed
permutation. A reflection with normal $\widehat n$ has the spin lift
$S_{\widehat n}=(\gamma\cdot\widehat n)\gamma_5$, for which
\[
 S_{\widehat n}(\gamma\cdot z)S_{\widehat n}^{-1}
       =\gamma\cdot(z-2(z\cdot\widehat n)\widehat n),
                         \qquad z\in\C^3.
\]
Choose a word representing $r$ in the three linear chamber
reflections and let $S_g$ be the corresponding product. It follows
that $S_g(\gamma\cdot z)S_g^{-1}=\gamma\cdot rz$.
The site map $(c,h)\mapsto(rc+v,rh)$ commutes with each bond
involution and carries $\widehat d_i(c,h)$ to
$r\widehat d_i(c,h)$. Therefore, for
\[
 (U_g\psi)(x)=S_g\psi(g^{-1}x),\qquad
 (g_*B)(c,c')=B(g^{-1}c,g^{-1}c'),
\]
the hopping formula gives
\begin{equation}\label{eq:response-complex-covariance}
 D^{\mathrm{ch}}[g_*B]+m
       =U_g(D^{\mathrm{ch}}[B]+m)U_g^{-1},
                                             \qquad B\ \text{complex}.
\end{equation}
This identity descends to the antiperiodic torus. Indeed
$r^{-1}e_\mu=\pm e_\nu$, and either translation by $ne_\nu$
or by $-ne_\nu$ changes an antiperiodic spinor's sign.
Thus $U_g$ and its inverse preserve the spinor space.
Taking determinants in \eqref{eq:response-complex-covariance}
shows that $N_0$, $c_j$, $a_{\mathrm{ch}}$ and $b_{\mathrm{ch}}$
are unchanged when the external cochain is replaced by $g_*B_t$.

For a physical translation $v\in\eta\Z^3$, cube averaging and
centered differences commute with translation. Periodization in
\eqref{eq:response-periodic-probe} then gives
$a_\eta[\beta(\,\cdot-v)](y)=a_\eta[\beta](y-v)$.
This is exactly the cochain action with lattice translation $v/\eta$.
The same determinant coefficients can therefore be used at each
lattice placement of the probe and at each of its reflected images.

\subsection{The two-photon coefficient}\label{subsec:response-polarization}

The coefficient is
\[
 a_{\mathrm{ch}}=\operatorname{Tr}(GD_2)-
                     \operatorname{Tr}(GD_1GD_1).
\]
Estimating the two traces separately would lose a power of the mesh.
We keep their cancellation by expressing their difference
as a polarization kernel. Its integral vanishes at zero momentum;
parity removes the first momentum derivative.

Define the bosonic momentum representatives
\[
 \mathcal K_{\eta,R}=
 \left\{\frac{2\pi}{R}j:j\in\Z^3,\ -n/2<j_\nu\le n/2\right\},
                   \qquad n=R/\eta,
\]
and the midpoint transform of the unperiodized probe
\begin{equation}\label{eq:response-midpoint-transform}
 h_{\eta,\nu}(K)=\eta^3
       \sum_{y\in\eta\Z^3+\eta e_\nu/2}
                     a_{\eta,\nu}^0(y)e^{-iK\cdot y}.
\end{equation}
The sum is finite. For real $p,k$ put
\begin{align}
 L_{\mu\nu}(p,k)
  &=\delta_{\mu\nu}\tr[R_m(p)W_\nu(p)]\notag\\
  &\quad-\tr[R_m(p)V_\mu(p+k/2)
                     R_m(p+k)V_\nu(p+k/2)],\label{eq:response-polarization-integrand}\\
 \Pi^{\eta,R}_{\mu\nu}(K)
  &=\eta^{-1}n^{-3}\sum_{p\in\mathcal P_n}
                       L_{\mu\nu}(p,\eta K),\notag\\
 \Pi^\eta_{\mu\nu}(K)
  &=\eta^{-1}\int_{[-\pi,\pi]^3}
          L_{\mu\nu}(p,\eta K)\frac{\dd p}{(2\pi)^3}.
                                  \label{eq:response-polarization-kernels}
\end{align}
Only $\Pi^{\eta,R}$ is a finite-volume polarization; the integral
$\Pi^\eta$ retains the lattice symbol but has no volume parameter.
The exact coefficient is
\begin{equation}\label{eq:response-polarization-identity}
 a_{\mathrm{ch}}(\eta,R)
    =R^{-3}\sum_{K\in\mathcal K_{\eta,R}}\sum_{\mu,\nu}
       \overline{h_{\eta,\mu}(K)}h_{\eta,\nu}(K)
                                      \Pi^{\eta,R}_{\mu\nu}(K).
\end{equation}

To check the midpoint and normalization in this formula, write
$b_\nu(c)=\eta A_\nu(c)$ and
$\widehat b_\nu(k)=n^{-3}\sum_cb_\nu(c)e^{-ik\cdot c}$.
The derivative vertices of \cref{subsec:response-contact}
give the blocks
\begin{align*}
 D_1[p+k,p]&=\sum_\nu\widehat b_\nu(k)e^{-ik_\nu/2}
                                          V_\nu(p+k/2),\\
 D_1[p,p+k]&=\sum_\mu
         \overline{\widehat b_\mu(k)e^{-ik_\mu/2}}
                                          V_\mu(p+k/2),\\
 D_2[p,p]&=\sum_\nu\sum_k|\widehat b_\nu(k)|^2W_\nu(p).
\end{align*}
Here $k$ ranges over $\eta\mathcal K_{\eta,R}$; the last
identity uses Parseval for the real function $b_\nu$.
For example, the positive hop in $D_1[p+k,p]$ is
$-i\widehat b_\nu(k)e^{ip_\nu}T_\nu^+$, and the negative hop
is $i\widehat b_\nu(k)e^{-i(p_\nu+k_\nu)}T_\nu^-$.
Their sum is the stated midpoint vertex. Under $k_\nu\mapsto
k_\nu+2\pi$ both the phase and that vertex change sign, so
their product is well-defined on momentum representatives.

Unfolding the periodization gives, for $k=\eta K$,
\[
 \widehat b_\nu(k)e^{-ik_\nu/2}
     =n^{-3}\eta\sum_{c\in\Z^3}
          a_{\eta,\nu}^0(\eta c+\eta e_\nu/2)
                            e^{-iK\cdot(\eta c+\eta e_\nu/2)}
     =\eta R^{-3}h_{\eta,\nu}(K).
\]
Substitute these three block identities into the two traces.
The scalar normalization is
$\eta^2R^{-6}=R^{-3}\eta^{-1}n^{-3}$, which proves
\eqref{eq:response-polarization-identity}. This identity
requires no separation of the probe from its periodic images.

\begin{proposition}\label{prop:response-polarization-bounds}
There is $C_\rho<\infty$, independent of $0<\bar m\le1$,
$0<\eta\le1$, $R/\eta\in\N$, $R/\eta\ge2$, and $K\in\R^3$,
such that
\begin{align}
 |\Pi^{\eta,R}_{\mu\nu}(K)-\Pi^\eta_{\mu\nu}(K)|
       &\le C_\rho\bar m^{-3}R^{-4},
                            \label{eq:response-polarization-quadrature}\\
 |\Pi^\eta_{\mu\nu}(K)|&\le C_\rho\bar m^{-1}|K|^2,\qquad
 |\nabla_K\Pi^\eta_{\mu\nu}(K)|
                    \le C_\rho\bar m^{-1}|K|.
                            \label{eq:response-polarization-quadratic}
\end{align}
\end{proposition}

\begin{proof}
We first estimate the quadrature error. If $F$ is a $C^4$
periodic function on the momentum torus, its Fourier coefficients
satisfy
\[
 |\widehat F(z)|\le |z|_\infty^{-4}
       \max_j\int_{[-\pi,\pi]^3}|\partial_j^4F(p)|
                                      \frac{\dd p}{(2\pi)^3},
                                                \qquad z\ne0.
\]
This follows by four integrations by parts in a coordinate attaining
$|z|_\infty$. The coefficients are absolutely summable, since
$\sum_{z\ne0}|z|_\infty^{-4}<42$. Their Fourier series equals
$F$: its difference from $F$ has zero Fourier coefficients,
and density of trigonometric polynomials in continuous periodic
functions makes the difference zero. The antiperiodic mean
selects the coefficients at $nz$ with sign $(-1)^{z_1+z_2+z_3}$.
Consequently
\begin{equation}\label{eq:response-quadrature-basic}
 \left|n^{-3}\sum_{p\in\mathcal P_n}F(p)
                  -\int_{[-\pi,\pi]^3}F(p)\frac{\dd p}{(2\pi)^3}\right|
 \le42n^{-4}\max_j
        \int_{[-\pi,\pi]^3}|\partial_j^4F(p)|
                                      \frac{\dd p}{(2\pi)^3}.
\end{equation}

Set $x(p)=\operatorname{dist}(p,2\pi\Z^3)+m$. On a period
cube $x(p)=|p|+m$ and $x(p)<7$. The inverse derivative
estimate \eqref{eq:response-inverse-derivatives}, used up
to order four, and the bounded symbol derivatives give
\begin{equation}\label{eq:response-polarization-fourth-derivative}
 |\partial_j^4L_{\mu\nu}(p,k)|
                 \le C_\rho[x(p)^{-6}+x(p+k)^{-6}].
\end{equation}
For the contact term, each Leibniz summand has at most five
powers of $x(p)^{-1}$; the remaining power is supplied by
$x(p)<7$. For the bubble term, a summand with $a$ derivatives
on the first inverse and $b$ on the second has the factor
$x(p)^{-1-a}x(p+k)^{-1-b}$, with $a+b\le4$.
The other derivatives fall on bounded vertices. Since
$u^{-s}v^{-t}\le u^{-s-t}+v^{-s-t}$ for $s,t>0$,
and $x<7$, the factor is bounded by a constant times the
right side of \eqref{eq:response-polarization-fourth-derivative}.
This bound is uniform in the shift $k$.

For $j>3$, periodicity and the radial integral give
\begin{equation}\label{eq:response-translated-integral}
 \int_{[-\pi,\pi]^3}x(p+u)^{-j}\,\dd p
 \le4\pi\int_0^\infty\frac{r^2\,\dd r}{(r+m)^j}
       =\frac{8\pi m^{3-j}}{(j-1)(j-2)(j-3)},
                                             \qquad u\in\R^3.
\end{equation}
Indeed the integral over a translated period cube is unchanged,
and on the original cube it is bounded by the integral over $\R^3$.
The case $j=6$ bounds the integral in
\eqref{eq:response-quadrature-basic} by $C_\rho m^{-3}$.
The definitions of the two polarizations then give the first
estimate, because
\[
 \eta^{-1}n^{-4}m^{-3}
           =\eta^{-1}(\eta/R)^4(\bar m\eta)^{-3}
           =\bar m^{-3}R^{-4}.
\]

For the quadratic bound, use the symmetric bubble
\[
 F_{\mu\nu}(p,k)=
 \tr[R_m(p-k/2)V_\mu(p)R_m(p+k/2)V_\nu(p)].
\]
It is periodic in $p$. Translating the period cube by $-k/2$
shows that its integral is the integral of the bubble in
\eqref{eq:response-polarization-integrand}. At $k=0$ the
whole integrand is a total derivative:
\begin{equation}\label{eq:response-polarization-total-derivative}
 L_{\mu\nu}(p,0)=
    \partial_\mu\tr[R_m(p)V_\nu(p)].
\end{equation}
In fact $\partial_\mu R_m=-R_mV_\mu R_m$ and
$\partial_\mu V_\nu=\delta_{\mu\nu}W_\nu$. Its integral
therefore vanishes, including the contact term. We obtain
\begin{equation}\label{eq:response-polarization-subtracted}
 \Pi^\eta_{\mu\nu}(K)=
     -\eta^{-1}\int_{[-\pi,\pi]^3}
          [F_{\mu\nu}(p,\eta K)-F_{\mu\nu}(p,0)]
                                      \frac{\dd p}{(2\pi)^3}.
\end{equation}

Conjugation by the involution in
\eqref{eq:response-bloch-parity} changes $(p,k)$ to $(-p,-k)$
in $F_{\mu\nu}$; the signs of its two vertices cancel.
Thus $F_{\mu\nu}(-p,-k)=F_{\mu\nu}(p,k)$.
Its first $k$ derivative at zero is odd in $p$, and hence
\begin{equation}\label{eq:response-polarization-parity-cancellation}
 \int_{[-\pi,\pi]^3}
              (\partial_{k_a}F_{\mu\nu})(p,0)\,\dd p=0.
\end{equation}
This removes the linear term as well as the constant term
removed by \eqref{eq:response-polarization-total-derivative}.

The Hessian has two derivatives in total on the inverse factors;
the midpoint vertices no longer depend on $k$. With $p_\pm=p\pm k/2$
and the vertices evaluated at $p$, it is
\begin{align*}
 \partial_{k_a}\partial_{k_b}F_{\mu\nu}
 =\tfrac14\big(&
  \tr[(\partial_a\partial_bR_m)(p_-)V_\mu R_m(p_+)V_\nu]\\
 &-\tr[(\partial_aR_m)(p_-)V_\mu(\partial_bR_m)(p_+)V_\nu]\\
 &-\tr[(\partial_bR_m)(p_-)V_\mu(\partial_aR_m)(p_+)V_\nu]\\
 &+\tr[R_m(p_-)V_\mu(\partial_a\partial_bR_m)(p_+)V_\nu]\big).
\end{align*}
Apply \eqref{eq:response-inverse-derivatives} through order two
and the same product inequality used above. It follows that
\[
 |\partial_{k_a}\partial_{k_b}F_{\mu\nu}(p,k)|
               \le C_\rho[x(p-k/2)^{-4}+x(p+k/2)^{-4}],
\]
and \eqref{eq:response-translated-integral} with $j=4$ gives
\begin{equation}\label{eq:response-polarization-hessian}
 \int_{[-\pi,\pi]^3}
     |\partial_{k_a}\partial_{k_b}F_{\mu\nu}(p,k)|
                             \frac{\dd p}{(2\pi)^3}
                  \le C_\rho m^{-1},\qquad k\in\R^3.
\end{equation}

Taylor's formula with integral remainder, followed by
\eqref{eq:response-polarization-parity-cancellation}, now makes
the two momentum powers explicit:
\begin{equation}\label{eq:response-polarization-taylor}
 \Pi^\eta_{\mu\nu}(K)=
   -\eta\int_{[-\pi,\pi]^3}\int_0^1(1-s)
      \sum_{a,b}K_aK_b
       (\partial_{k_a}\partial_{k_b}F_{\mu\nu})(p,s\eta K)
                           \,\dd s\,\frac{\dd p}{(2\pi)^3}.
\end{equation}
This gives the bound $C_\rho(\eta/m)|K|^2$.
For the gradient, differentiate
\eqref{eq:response-polarization-subtracted}, subtract its
zero-momentum derivative using parity, and apply the
first-order integral remainder:
\[
 \partial_{K_a}\Pi^\eta_{\mu\nu}(K)=
  -\eta\int_{[-\pi,\pi]^3}\int_0^1
       \sum_bK_b(\partial_{k_a}\partial_{k_b}F_{\mu\nu})(p,s\eta K)
                              \,\dd s\,\frac{\dd p}{(2\pi)^3}.
\]
All differentiations are under an integral over a compact cube
of a smooth function, since $m>0$. Using
\eqref{eq:response-polarization-hessian} and $\eta/m=\bar m^{-1}$
proves \eqref{eq:response-polarization-quadratic}.
\end{proof}

\subsection{Uniform bounds on the determinant coefficients}
\label{subsec:response-uniform-coefficients}

We finish by pairing the polarization bounds with decay of the
probe transform. This step needs no continuum evaluation of
the two-photon kernel.

For each integer $j\ge0$ there is $C_{\beta,j}<\infty$ such that
\begin{equation}\label{eq:response-midpoint-decay}
 |h_{\eta,\nu}(K)|\le C_{\beta,j}(1+|K|)^{-j},
                \qquad |K|_\infty\le\pi/\eta,\quad0<\eta\le1.
\end{equation}
The constants are unchanged by any real translation of $\beta$.
Here is the discrete integration by parts that makes this
uniform in the mesh. Suppose a smooth function $v$ is supported
in a cube of side $\ell$, and sum it over a translate of $\eta\Z^3$.
The backward difference
$\Delta_{\eta,\mu}v(y)=\eta^{-1}(v(y)-v(y-\eta e_\mu))$
satisfies
$\|\Delta_{\eta,\mu}^jv\|_\infty
\le\|\partial_\mu^jv\|_\infty$, by its integral representation.
Its support is enlarged by at most $j\eta$ in one direction,
and the number of sampling points, multiplied by $\eta^3$,
is at most $(\ell+j+1)^3$. Summation by parts gives
\[
 [\eta^{-1}(1-e^{-i\eta K_\mu})]^j
       \eta^3\sum_yv(y)e^{-iK\cdot y}
          =\eta^3\sum_y\Delta_{\eta,\mu}^jv(y)e^{-iK\cdot y}.
\]
On $|\eta K_\mu|\le3\pi/2$ the modulus of the bracket
is at least $0.3|K_\mu|$. Choose a coordinate attaining
$|K|_\infty$. For $|K|\ge1$ this is at least
$(1+|K|)/12$; for $|K|<1$ use the point-count estimate directly.
We thus obtain, on that momentum cube,
\[
 \left|\eta^3\sum_yv(y)e^{-iK\cdot y}\right|
 \le12^j(\ell+j+1)^3
       \max(\|v\|_\infty,\max_\mu\|\partial_\mu^jv\|_\infty)
                                               (1+|K|)^{-j}.
\]

Apply this to $v=a_{\eta,\nu}^0$ on
$\eta\Z^3+\eta e_\nu/2$. Cube averaging has norm one on
derivative sup norms, and a centered difference is an average
of the corresponding first derivative. Consequently
$\|\partial_\mu^ja_{\eta,\nu}^0\|_\infty
\le2|\beta|_{C^{j+1}}$, while its support has side at most
$D_\beta+2\eta$. One may take
\[
 C_{\beta,j}=2\cdot12^j(D_\beta+j+3)^3|\beta|_{C^{j+1}}.
\]
This proves \eqref{eq:response-midpoint-decay}, with position-independent
constants and no assumption on $R$.

\begin{proposition}\label{prop:response-uniform-coefficients}
For each fixed $\rho$, $\bar m$ and $\beta$ as above,
\[
 |c_2|\le C_{\rho,\bar m,\beta},
              \qquad 0<\eta\le1,\quad R\ge1,\quad
                         R/\eta\in\N,\quad R/\eta\ge2.
\]
If in addition $R\ge D_\beta+1$, then
$|c_4|\le C_{\rho,\bar m,\beta}$.
The bounds are uniform in all real translations of $\beta$,
and the coefficients have no dependence on $e$.
\end{proposition}

\begin{proof}
\Cref{prop:response-polarization-bounds} gives
\[
 |\Pi^{\eta,R}_{\mu\nu}(K)|
   \le C_\rho(\bar m^{-3}R^{-4}+\bar m^{-1}|K|^2).
\]
In \eqref{eq:response-polarization-identity} use the decay
estimate with $j=2$ for the quadrature term and $j=3$ for
the quadratic term. Since $R\ge1$, both contributions are
bounded by a constant times $(1+|K|)^{-4}$; hence
\[
 |c_2|=|a_{\mathrm{ch}}|
 \le9C_\rho
       (\bar m^{-3}C_{\beta,2}^2+\bar m^{-1}C_{\beta,3}^2)
       R^{-3}\sum_{K\in\mathcal K_{\eta,R}}(1+|K|)^{-4}.
\]
The last sum is bounded uniformly for $R\ge1$. To see this,
tile $\R^3$ with cubes of side $2\pi/R$ centered at
$K\in(2\pi/R)\Z^3$. On each cube,
$1+|P|\le(1+\sqrt3\pi)(1+|K|)$, so its volume times
$(1+|K|)^{-4}$ is bounded by a constant times the integral
of $(1+|P|)^{-4}$ there. Sum over the cubes; that integral
over $\R^3$ is finite. Restricting the sum to
$\mathcal K_{\eta,R}$ only decreases it. This proves the
uniform second-order bound.

For the fourth coefficient, use
\eqref{eq:response-coefficient-relations} and
the fourth-trace estimates \eqref{eq:response-fourth-errors}, proved in \cref{subsec:response-fourth}. The infinite-space
trace $b(f)$ is finite by \eqref{eq:response-four-ideal-infinite},
and the two errors in \eqref{eq:response-fourth-errors} are
bounded for $R\ge D_\beta+1$ and $0<\eta\le1$.
Thus $|b_{\mathrm{ch}}|$ is bounded, and adding $3c_2^2$
gives the asserted bound for $c_4$. All constants in the
proof depend on support sizes and derivative norms of $\beta$,
not on its position. Independence of $e$ follows directly
from the external connection convention.
\end{proof}



\begin{proof}[Completion of the proof of \cref{prop:response-fourth}]
The fourth-trace assertions were proved in
\Cref{subsec:response-fourth}.
The determinant symmetries and coefficient identities follow from
\Cref{subsec:response-symmetry}, and
\Cref{prop:response-uniform-coefficients} gives the second- and
fourth-order bounds with the stated volume ranges.
\end{proof}

\printbibliography
\end{document}
